\documentclass[11pt]{article}

\usepackage[margin=1in]{geometry}
\usepackage{amsmath,amssymb,amsthm}
\usepackage[hidelinks]{hyperref}
\usepackage{enumitem}

\setlist[itemize]{leftmargin=1.5em}
\setlength{\parindent}{0pt}
\setlength{\parskip}{0.7em}

\newtheorem*{claim}{Claim}

\begin{document}

\begin{center}
{\Large Literature-Already-Answered Packet}\\[0.25em]
{\large Upper Khintchine estimate in noncommutative symmetric spaces}\\[0.35em]
\textbf{Status:} literature already answered, not a new automatic-discovery proof
\end{center}

\section*{Source and later answer}

The source paper is:
\[
\text{arXiv:0803.4404, Christian Le Merdy and Fedor Sukochev,}
\]
\[
\text{\emph{Rademacher averages on noncommutative symmetric spaces}, 2008.}
\]
It appeared in \emph{Journal of Functional Analysis} 255 (2008),
3329--3355. The local copy is included as \texttt{source\_paper.pdf}.

The later supporting paper is:
\[
\text{S. Dirksen, B. de Pagter, D. Potapov, and F. Sukochev,}
\]
\[
\text{\emph{Rosenthal inequalities in noncommutative symmetric spaces},}
\]
\[
\text{\emph{Journal of Functional Analysis} 261 (2011), no. 10, 2890--2925,}
\]
\[
\text{DOI: \href{https://doi.org/10.1016/j.jfa.2011.07.015}{10.1016/j.jfa.2011.07.015}.}
\]

\section*{Open question in the source}

Let \(E\) be a separable, or dual-separable, symmetric function space, let
\(M\) be a semifinite von Neumann algebra, and let \(E(M)\) be the associated
noncommutative symmetric space. For a finite family \((x_k)\subset E(M)\), the
source defines
\[
\|(x_k)\|_{\max}
= \max\left\{
\left\|\left(\sum_k x_k^*x_k\right)^{1/2}\right\|_{E(M)},
\left\|\left(\sum_k x_kx_k^*\right)^{1/2}\right\|_{E(M)}
\right\}.
\]
Theorem 1.1(2) of the source proves, under \(q_E<\infty\) and \(p_E>1\), the
upper Rademacher estimate
\[
\left\|\sum_k \varepsilon_k\otimes x_k\right\|_E
\lesssim_E \|(x_k)\|_{\max}.
\tag{1}
\]
Immediately after Theorem 1.1, in the introduction, the authors state as an
open problem that they do not know whether this second part remains true
without assuming \(p_E>1\).

\section*{Supporting answer}

The abstract of Dirksen--de Pagter--Potapov--Sukochev (2011) states that the
paper gives a direct proof of the upper Khintchine inequality for a
noncommutative symmetric quasi-Banach function space with nontrivial upper
Boyd index, and that this settles an open question of Le Merdy and Sukochev
(2008).

The condition ``nontrivial upper Boyd index'' is the condition \(q_E<\infty\)
in the notation of the source paper. Thus the supporting result proves the
upper Khintchine estimate (1) under the source paper's \(q_E<\infty\)
hypothesis alone, removing the extra assumption \(p_E>1\).

\begin{claim}
The open problem after Theorem 1.1 in arXiv:0803.4404 is already answered by
Dirksen--de Pagter--Potapov--Sukochev, \emph{Rosenthal inequalities in
noncommutative symmetric spaces}, JFA 261 (2011), 2890--2925.
\end{claim}

\begin{proof}
This is a direct literature identification. The source question asks exactly
whether the upper Khintchine/Rademacher estimate in Theorem 1.1(2) remains
true when the assumption \(p_E>1\) is removed while retaining \(q_E<\infty\).
The supporting paper explicitly says that it proves the upper Khintchine
inequality for noncommutative symmetric quasi-Banach function spaces with
nontrivial upper Boyd index and that this settles the Le Merdy--Sukochev
open question. Since nontrivial upper Boyd index is \(q_E<\infty\), this is
the requested removal of \(p_E>1\). Therefore the extracted open problem has
a separate later literature answer.
\end{proof}

\section*{Classification and scope}

This packet belongs in
\texttt{solutions/literature\_already\_answered/}. The supporting authors
explicitly identify their result as settling the open question from Le
Merdy--Sukochev (2008). This is a retrieval/status result, not a new proof by
the run.

The recorded answer is limited to the upper Khintchine estimate in Theorem
1.1(2). The source paper also discusses examples, optimality issues, and
double sums; this packet makes no additional claim about those passages unless
checked separately.

\section*{Verification notes}

\begin{itemize}
\item The source is not a same-paper false positive: the open problem appears
  immediately after Theorem 1.1 in the 2008 paper.
\item The supporting source is separate and later, and its abstract explicitly
  says that it settles the Le Merdy--Sukochev open question.
\item The decisive match is between \(q_E<\infty\) in the source theorem and
  nontrivial upper Boyd index in the supporting paper.
\end{itemize}

\section*{References}

\begin{itemize}
\item C. Le Merdy and F. Sukochev, arXiv:0803.4404,
  \emph{Rademacher averages on noncommutative symmetric spaces},
  J. Funct. Anal. 255 (2008), 3329--3355.
\item S. Dirksen, B. de Pagter, D. Potapov, and F. Sukochev,
  \emph{Rosenthal inequalities in noncommutative symmetric spaces},
  J. Funct. Anal. 261 (2011), no. 10, 2890--2925.\\
  DOI: \href{https://doi.org/10.1016/j.jfa.2011.07.015}{10.1016/j.jfa.2011.07.015}.
\end{itemize}

\end{document}
